\chapter{Glossary}

\begin{description}
\item[Arbitrage] A strategy that has no initial net cost, no possible loss, and a positive gain in at least one state under stated assumptions.
\item[Asset] A controlled resource expected to provide future economic benefits.
\item[Basis] A difference between related spot and derivative prices; the sign convention must be stated.
\item[Beta] Covariance of an asset's return with a factor divided by factor variance.
\item[Bond] A debt claim promising contractual payments and principal under stated terms.
\item[Call option] The right, not obligation, to buy an underlying at a specified strike.
\item[Capital structure] The mix of debt, equity, and other financing claims.
\item[Clearing] Processes that calculate obligations, manage collateral, and support settlement.
\item[Convexity] Curvature of a price-yield or value-risk relationship.
\item[Correlation] Standardised linear dependence between two variables.
\item[Counterparty risk] Risk that the other party fails to perform.
\item[Coupon] Contractual interest payment on a bond.
\item[Credit spread] Yield or price compensation relative to a benchmark claim; not pure default probability.
\item[Derivative] A contract whose value depends on an underlying asset, rate, index, or event.
\item[Discount factor] Present value of one unit paid at a future date.
\item[Discount rate] Rate used to translate future cash flows into present value.
\item[Duration] A time-weighted or sensitivity-based measure of bond price exposure to yield.
\item[Efficient frontier] Feasible portfolios that are not dominated in expected return and risk under a chosen model.
\item[Equity] Residual ownership claim after liabilities.
\item[Expected shortfall] Average tail loss beyond a chosen loss quantile, under a defined distribution.
\item[Factor] A systematic return or state variable used to explain exposures.
\item[Forward] A private agreement to transact in the future at a fixed delivery price.
\item[Futures] A standardised, margined derivative usually marked to market daily.
\item[Gamma] Second derivative of derivative value with respect to underlying price.
\item[Hedging] Taking an offsetting position to reduce a defined exposure.
\item[Implied volatility] Volatility parameter that makes a model price equal an observed option price.
\item[Leverage] Exposure to assets or risks relative to equity or capital.
\item[Liquidity] Ability to transact or meet obligations with limited price impact or delay.
\item[Market capitalisation] Share price multiplied by shares outstanding.
\item[Market microstructure] Study of trading mechanisms, order flow, spreads, and price formation.
\item[Monte Carlo] Estimation by averaging outcomes from simulated random paths.
\item[Net asset value] Assets minus liabilities divided by units or shares of a pooled vehicle.
\item[No-arbitrage] Pricing principle that rules out exploitable riskless discrepancies under a model.
\item[Nominal return] Return measured in money prices before removing inflation.
\item[Option] A contingent right with nonlinear payoff.
\item[Portfolio] A collection of financial claims and their weights.
\item[Present value] Current equivalent value of future cash flows under a discounting convention.
\item[Put option] The right, not obligation, to sell an underlying at a specified strike.
\item[Risk premium] Expected compensation beyond a benchmark for bearing a risk or providing a service.
\item[Risk-neutral measure] A pricing probability measure under which appropriately discounted tradable prices are martingales.
\item[Sharpe ratio] Mean excess return divided by the standard deviation of excess return.
\item[Spot rate] Rate used to discount a single cash flow to a specific maturity.
\item[Stress test] Analysis of outcomes under a specified adverse scenario.
\item[Swap] Contract exchanging cash-flow streams, often based on rates, currencies, or commodities.
\item[Systemic risk] Risk that distress spreads through common exposures, interconnections, or feedback loops.
\item[Terminal value] Value assigned to cash flows beyond an explicit forecast horizon.
\item[Theta] Sensitivity of option value to the passage of time under a convention.
\item[Value at Risk] A loss quantile at a stated horizon and confidence level.
\item[Vega] Sensitivity of option value to the volatility input.
\item[Volatility] A measure of return dispersion, commonly standard deviation under a chosen horizon.
\end{description}
